% GENERATED FILE. Edit canonical lessons, then run npm run generate:books.
% canonical-source-hash: fnv1a64:58c911d503cc72b5
% canonical-lessons: GU-C271-namra, GU-C271-udar, GU-C271-svarthi, GU-C271-murkh, GU-C271-jiddi

\chapter{Polite or selfish}
\label{ch:gu-polite-or-selfish}
\hlchaptermodality{271}

\begin{chapteropening}
\textbf{By the end of this chapter:} \emph{I can describe people as polite, generous or stubborn and ask whether they are.}

Complete the last lesson of chapter 271: I can describe people as polite, generous or stubborn and ask whether they are.
\end{chapteropening}

\section[\texorpdfstring{\gu{નમ્ર} (namra)}{નમ્ર (namra)}]{\gu{નમ્ર} (namra) — polite, humble}

\label{lesson:GU-C271-namra}



\begin{quote}
Before the new one: say the Gujarati for funny, then the Gujarati for arrogant.
\end{quote}

\subsection*{You'll want to know: \gu{નમ્ર}}
\textbf{\gu{નમ્ર}} — \emph{namra} — \textquotedblleft{}polite, humble\textquotedblright{}.

\begin{quote}
\textbf{\gu{શું} \gu{તમારો} \gu{દીકરો} \gu{નમ્ર} \gu{છે}?} — \emph{śũ tamāro dīkro namra che?} — \textquotedblleft{}Is your son polite?\textquotedblright{}
\end{quote}

\begin{cousinweb}
\textbf{\gu{નમ્ર}} (\emph{namra}) is Sanskrit \emph{namrá}, \textquotedblleft{}bent, submissive\textquotedblright{}, taken over whole as a tatsama (CDIAL 6964).
\end{cousinweb}

\begin{grammarlens}[title={What you've built}]
One more word for everyday life at A2.
\end{grammarlens}

\subsection*{Your turn}
\begin{itemize}
\raggedright
  \item \emph{Say it:} \emph{namra}
  \item \emph{Say it:} \emph{namra}, once more
  \item \emph{Say it:} \emph{śũ tamāro dīkro namra che?} — Is your son polite?
  \item \emph{Recall:} say the Gujarati for funny, then the Gujarati for arrogant, then say \emph{namra} again
\end{itemize}

\begin{tcolorbox}[breakable,colback=teal!4,colframe=teal!35!black,title={Before you move on}]
What does \gu{નમ્ર} mean? (\textquotedblleft{}Polite, humble\textquotedblright{}.) Say it once more.
\end{tcolorbox}
\section[\texorpdfstring{\gu{ઉદાર} (udār)}{ઉદાર (udār)}]{\gu{ઉદાર} (udār) — generous}

\label{lesson:GU-C271-udar}



\begin{quote}
Before the new one: say the Gujarati for arrogant, then the Gujarati for polite.
\end{quote}

\subsection*{You'll want to know: \gu{ઉદાર}}
\textbf{\gu{ઉદાર}} — \emph{udār} — \textquotedblleft{}generous\textquotedblright{}.

\begin{quote}
\textbf{\gu{અમારા} \gu{પડોશી} \gu{ઉદાર} \gu{નથી}.} — \emph{amārā paḍośī udār nathī} — \textquotedblleft{}Our neighbours are not generous.\textquotedblright{}
\end{quote}

\begin{cousinweb}
\textbf{\gu{ઉદાર}} (\emph{udār}) is Sanskrit \emph{udārá}, \textquotedblleft{}exciting\textquotedblright{} and later \textquotedblleft{}lofty, noble\textquotedblright{}, taken over as a tatsama (CDIAL 1935).
\end{cousinweb}

\begin{grammarlens}[title={What you've built}]
One more word for everyday life at A2.
\end{grammarlens}

\subsection*{Your turn}
\begin{itemize}
\raggedright
  \item \emph{Say it:} \emph{udār}
  \item \emph{Say it:} \emph{udār}, once more
  \item \emph{Say it:} \emph{amārā paḍośī udār nathī} — Our neighbours are not generous.
  \item \emph{Recall:} say the Gujarati for arrogant, then the Gujarati for polite, then say \emph{udār} again
\end{itemize}

\begin{tcolorbox}[breakable,colback=teal!4,colframe=teal!35!black,title={Before you move on}]
What does \gu{ઉદાર} mean? (\textquotedblleft{}Generous\textquotedblright{}.) Say it once more.
\end{tcolorbox}
\section[\texorpdfstring{\gu{સ્વાર્થી} (svārthī)}{સ્વાર્થી (svārthī)}]{\gu{સ્વાર્થી} (svārthī) — selfish}

\label{lesson:GU-C271-svarthi}



\begin{quote}
Before the new one: say the Gujarati for polite, then the Gujarati for generous.
\end{quote}

\subsection*{You'll want to know: \gu{સ્વાર્થી}}
\textbf{\gu{સ્વાર્થી}} — \emph{svārthī} — \textquotedblleft{}selfish\textquotedblright{}.

\begin{quote}
\textbf{\gu{મારો} \gu{દોસ્ત} \gu{સ્વાર્થી} \gu{નથી}.} — \emph{māro dost svārthī nathī} — \textquotedblleft{}My pal is not selfish.\textquotedblright{}
\end{quote}

\begin{cousinweb}
\textbf{\gu{સ્વાર્થી}} (\emph{svārthī}) is built on Sanskrit \emph{svārtha}, \textquotedblleft{}one's own aim\textquotedblright{}, from \emph{sva}, \textquotedblleft{}own\textquotedblright{}, and \emph{ártha}, \textquotedblleft{}aim\textquotedblright{} (CDIAL 638), with \emph{-ī}; a learned word, a tatsama in shape.
\end{cousinweb}

\begin{grammarlens}[title={What you've built}]
One more word for everyday life at A2.
\end{grammarlens}

\subsection*{Your turn}
\begin{itemize}
\raggedright
  \item \emph{Say it:} \emph{svārthī}
  \item \emph{Say it:} \emph{svārthī}, once more
  \item \emph{Say it:} \emph{māro dost svārthī nathī} — My pal is not selfish.
  \item \emph{Recall:} say the Gujarati for polite, then the Gujarati for generous, then say \emph{svārthī} again
\end{itemize}

\begin{tcolorbox}[breakable,colback=teal!4,colframe=teal!35!black,title={Before you move on}]
What does \gu{સ્વાર્થી} mean? (\textquotedblleft{}Selfish\textquotedblright{}.) Say it once more.
\end{tcolorbox}
\section[\texorpdfstring{\gu{મૂર્ખ} (mūrkh)}{મૂર્ખ (mūrkh)}]{\gu{મૂર્ખ} (mūrkh) — foolish, a fool}

\label{lesson:GU-C271-murkh}



\begin{quote}
Before the new one: say the Gujarati for generous, then the Gujarati for selfish.
\end{quote}

\subsection*{You'll want to know: \gu{મૂર્ખ}}
\textbf{\gu{મૂર્ખ}} — \emph{mūrkh} — \textquotedblleft{}foolish, a fool\textquotedblright{}.

\begin{quote}
\textbf{\gu{શું} \gu{હું} \gu{મૂર્ખ} \gu{છું}?} — \emph{śũ hũ mūrkh chũ?} — \textquotedblleft{}Am I a fool?\textquotedblright{}
\end{quote}

\begin{cousinweb}
\textbf{\gu{મૂર્ખ}} (\emph{mūrkh}) is Sanskrit \emph{mūrkha}, \textquotedblleft{}foolish, stupid\textquotedblright{}, taken over whole as a tatsama. CDIAL has no entry for it.
\end{cousinweb}

\begin{grammarlens}[title={What you've built}]
One more word for everyday life at A2.
\end{grammarlens}

\subsection*{Your turn}
\begin{itemize}
\raggedright
  \item \emph{Say it:} \emph{mūrkh}
  \item \emph{Say it:} \emph{mūrkh}, once more
  \item \emph{Say it:} \emph{śũ hũ mūrkh chũ?} — Am I a fool?
  \item \emph{Recall:} say the Gujarati for generous, then the Gujarati for selfish, then say \emph{mūrkh} again
\end{itemize}

\begin{tcolorbox}[breakable,colback=teal!4,colframe=teal!35!black,title={Before you move on}]
What does \gu{મૂર્ખ} mean? (\textquotedblleft{}Foolish, a fool\textquotedblright{}.) Say it once more.
\end{tcolorbox}
\section[\texorpdfstring{\gu{જિદ્દી} (jiddī)}{જિદ્દી (jiddī)}]{\gu{જિદ્દી} (jiddī) — stubborn}

\label{lesson:GU-C271-jiddi}



\begin{quote}
Before the new one: say the Gujarati for selfish, then the Gujarati for foolish.
\end{quote}

\subsection*{You'll want to know: \gu{જિદ્દી}}
\textbf{\gu{જિદ્દી}} — \emph{jiddī} — \textquotedblleft{}stubborn\textquotedblright{}.

\begin{quote}
\textbf{\gu{મારો} \gu{ભત્રીજો} \gu{જિદ્દી} \gu{છે}, \gu{તે} \gu{સાંભળતો} \gu{નથી}.} — \emph{māro bhatrījo jiddī che, te sā̃bhaḷto nathī} — \textquotedblleft{}My nephew is stubborn; he does not listen.\textquotedblright{}
\end{quote}

\begin{cousinweb}
\textbf{\gu{જિદ્દી}} (\emph{jiddī}) is built on \emph{jid}, \textquotedblleft{}obstinacy\textquotedblright{}, with \emph{-ī}. \emph{Jid} is usually traced to Arabic \emph{ḍidd}, \textquotedblleft{}an opposite\textquotedblright{}, through Persian; Lane gives the root \emph{ḍ-d-d}, \textquotedblleft{}to overcome in contention\textquotedblright{} (Lane).
\end{cousinweb}

\begin{grammarlens}[title={What you've built}]
One more word for everyday life at A2.
\end{grammarlens}

\subsection*{Your turn}
\begin{itemize}
\raggedright
  \item \emph{Say it:} \emph{jiddī}
  \item \emph{Say it:} \emph{jiddī}, once more
  \item \emph{Say it:} \emph{māro bhatrījo jiddī che, te sā̃bhaḷto nathī} — My nephew is stubborn; he does not listen.
  \item \emph{Recall:} say the Gujarati for selfish, then the Gujarati for foolish, then say \emph{jiddī} again
\end{itemize}

\begin{tcolorbox}[breakable,colback=teal!4,colframe=teal!35!black,title={Before you move on}]
What does \gu{જિદ્દી} mean? (\textquotedblleft{}Stubborn\textquotedblright{}.) Say it once more.
\end{tcolorbox}
